\documentclass[10pt,a4paper]{amsart}
\usepackage[margin=19mm]{geometry}
\usepackage{amsmath,amssymb,mathtools}
\usepackage[scr=rsfs,cal=euler]{mathalfa}
\usepackage{xfrac}
\usepackage[unicode,hidelinks,bookmarks=false]{hyperref}
\newcommand{\ie}{i.e.\ }
\newcommand{\cf}{cf.\ }
\newcommand{\resp}{respectively\ }
\newcommand{\Subsection}{\subsection}
\providecommand{\nobreakdash}{\nobreak\hbox{-}}
\setlength{\emergencystretch}{2em}
\multlinegap0pt
\usepackage[normalem]{ulem}
\usepackage{amssymb}
\usepackage{tikz}
\newtheorem{lemma}{Lemma}
\newcommand{\Iso}{\textup{Iso}}
\renewcommand{\L}{\Lambda}
\renewcommand{\a}{\alpha}
\renewcommand{\b}{\beta}
\newcommand{\de}{\delta}
\newcommand{\ga}{\gamma}
\title{Uniqueness theorems for combinatorial C*-algebras}
\author{Charles Starling}
\date{Source: arXiv:2604.21116v2 (CC BY 4.0)}
\begin{document}
\maketitle

\noindent\emph{Source and licence.} Uniqueness theorems for combinatorial C*-algebras. Version arXiv:2604.21116v2 (CC BY 4.0), distributed by arXiv under the Creative Commons Attribution 4.0 International licence, http://creativecommons.org/licenses/by/4.0/, as recorded in the arXiv licence metadata for this exact version and shown on the abstract page https://arxiv.org/abs/2604.21116v2. Full licence notice: \texttt{licenses/arxiv-2604.21116-CC-BY-4.0.txt}.\par\smallskip

\noindent\textit{Editorial scope.} Counted unit: the article's lemma lemma (source0.tex lines 572--574). The article's complete original proof of it is delivered with it (source0.tex lines 575--609, one \texttt{proof} environment of 35 lines). No mathematics is written, repaired or re-derived: the statement and the proof are the article's own lines, in the article's own notation, and no retained line is substituted or re-flowed.  Retained uncounted as the source-local context the proof needs: lines 545--548.  Nothing was omitted from any retained span, so the package carries no omission record.  No retained line contains a citation, so the wrapper carries no bibliography and no bibliography entry.  No retained line contains a figure environment, an image-inclusion command or drawing code, so no image is delivered and the package declares no asset dependency.  Editorial formatting only, in addition: an AMS \texttt{amsart} preamble that restates the article's own declarations and macros used by the retained lines, namely the environments \texttt{lemma} on separate counters, together with the standard packages those lines need.  The retained environments print in this document as Lemma 1 in order of appearance, whereas the article prints the same items with the numbering of its own full document.  The complete native source of the article as distributed in the arXiv e-print for this exact version is bundled unchanged as \texttt{sources/199000/source0.tex}.
	\begin{lemma}\label{lem:iso_intersection}
		Let $\L$ be a finitely aligned LCSC, and suppose that $\alpha\beta^*\in S_\L^\Iso$ with $\alpha\L\cap \beta\L = \bigcup_{i=1}^n\gamma_i\L$. Then $D_{\alpha\L} = D_{\beta\L} = \bigcup_{i=1}^nD_{\gamma_i\L}$. Furthermore, if $\gamma\L\cap \alpha\L \neq \emptyset$ or $\gamma\L\cap \beta\L \neq \emptyset$, then there exists $1\leq i\leq n$ such that $\gamma\L\cap\gamma_i\L\neq \emptyset$.
		
	\end{lemma}

	\begin{lemma}
		$F_\L$ is an inverse subsemigroup of $S_\L^\Iso$. 
	\end{lemma}

	\begin{proof}
		Clearly $F_\L$ is closed under inverses, so we will be done if we can show it is closed under products. Take $s,t\in F_\L$ and write
		\[
		s = \bigcup_{n=1}^N \a_n\b_n^*, \hspace{1cm} t = \bigcup_{m=1}^M \ga_m\de_m^*.
		\]
		The product will be zero unless there is $x\in \L$ such that $\a_n, \b_n, \ga_m, \de_m\in x\L$ for all $1\leq n\leq N$ and $1\leq m\leq M$. In the case of a nonzero product, exhaustiveness of $\{\ga_m\}$ and $\{\b_n\}$ imply that for each $1\leq n\leq N$ and $1\leq m\leq M$ there exists $K(n,m)\geq 1$ and $\tau_k^{(n,m)}\in \Lambda$ for each $1\leq k\leq K(n,m)$ such that
		\begin{equation}\label{eq:bn_gm}
		\b_n \L \cap \ga_m\L = \bigcup_{k=1}^{K(n,m)} \tau_k^{(n,m)}\L
		\end{equation}
		so that, for each  $1\leq k\leq K(n,m)$ we have
		\begin{equation}\label{eq:tau_b_ga}
		\tau_k^{(n,m)} = \b_n\b_{n,m,k} = \ga_m\ga_{n,m,k}
		\end{equation}
		for elements $\b_{n,m,k}, \ga_{n,m,k}\in \L$. Then one has (by distributativity in the symmetric inverse monoid)
		\[
		\a_n\b_n^*\ga_m\de^*_n = \a_n\left(\bigcup_{k=1}^{K(n,m)}\b_{n,m,k}\ga_{n,m,k}^*\right)\de_m^* = \bigcup_{k=1}^{K(n,m)}\a_n\b_{n,m,k}(\de_m\ga_{n,m,k})^* .
		\]
		So again by distributativity in the symmetric inverse monoid we have
		\[
		st = \bigcup_{n=1}^N\bigcup_{m=1}^M\bigcup_{k=1}^{K(n,m)}\a_n\b_{n,m,k}(\de_m\ga_{n,m,k})^*.
		\]
		We will be done if we can show that $\{\a_n\b_{n,m,k}\}$ and $\{\de_m\ga_{n,m,k}\}$ are exhaustive (at $x$).
		
		Let $\eta\in x\L$ and find $\b_n$ such that $\b_n\L\cap \eta\L\neq\emptyset$. By Lemma~\ref{lem:iso_intersection} we can find $\rho\in \a_n\L\cap\b_n\L$ such that $\eta\L\cap \rho\L\neq \emptyset$. Write $\rho = \a_n\a' = \b_n\b'$ for $\a',\b'\in\Lambda$ and find $\eta', \rho''\in\Lambda$ such that $\eta\eta' = \rho\rho''$. 
		
		The element $\b_n\alpha'\rho''$ is well-defined and has range $x$, so we can find $m$ such that $\ga_m\L\cap \b_n\alpha'\rho''\L\neq\emptyset$. If $\zeta\in \ga_m\L\cap \b_n\alpha'\rho''\L\neq\emptyset$ we can, by \eqref{eq:bn_gm} find $\tau', \ga',\b''\in\L$ and $1\leq k\leq K(n,m)$ such that
		\[
		\zeta = \ga_m\ga' = \b_n\alpha'\rho''\b'' = \tau_k^{(n,m)}\tau' \stackrel{\eqref{eq:tau_b_ga}}{=} \b_n\b_{n,m,k}\tau'.
		\]
		Thus left cancellativity implies $\b_{n,m,k}\tau' = \alpha'\rho''\b''$. And so we have
		\[
		\a_n\b_{n,m,k}\tau' = \a_n\alpha'\rho''\b'' = \rho\rho''\b'' = \eta\eta'
		\]
		implying that $\eta\Lambda\cap \a_n\b_{n,m,k}\Lambda\neq\emptyset$. Since $\eta\in x\L$ was arbitrary, $\{\a_n\b_{n,m,k}\}$ is exhaustive. Applying inverses and running the same argument gives that $\{\de_m\ga_{n,m,k}\}$ is also exhaustive, and so $F_\L$ is closed under products.
	\end{proof}

\end{document}
